% Terminal note of the homily: everything the speech itself must not carry ---
% the audience, the spoken word count and the pace it implies, the relation to
% the three interpretations of the expansive study, the connections that are
% the homily's own, the translation each quoted text is given in, and the exact
% loci --- followed by the References. Nothing here is read aloud.

\section*{Note on Sources and Delivery}
\runninghead{Sources and Delivery}

{\footnotesize

\textbf{Audience and occasion.} An adult parish assembly, at the Mass
\latin{Dominica sexta quae superfuit post Epiphaniam}, II classis, of the
1962 \work{Missale Romanum}, editio typica, pp.~421--422, nos.~1759--1768,
said on Sunday 15~November 2026, the twenty-fifth Sunday after Pentecost. The
speech presumes that it follows the Gospel, and it speaks of the Offertory,
Communion and Postcommunion as still to come. It does not use the Introit,
the Gradual, the Secret or the Preface.

\textbf{Length and pace.} The spoken body runs to about 1,300 words (1,284 by
a plain count with comments and markup removed), which is about 9.9 to 10.7
minutes at an unhurried 120 to 130 words a minute, by arithmetic alone; no
one has delivered or timed it. It was read through in full, silently, for
sense and ease of speech.

\textbf{Relation to the reviewed interpretations.} The homily follows the
second of the expansive study's three interpretations, \textit{Leaven Hidden
in Three Measures: Love at Work in the Whole Person}. It adds from the first
interpretation only Bede's sentence on the mustard seed, that the preaching
grows in the mind of each believer because no one is made perfect all at
once, which that interpretation cites at the Postcommunion. The other
answers to what the leaven is are left to the studies: Jerome's knowledge of
the Scriptures (\work{In Matthaeum} II, on 13:33, PL~26 cols.~91--92),
Hilary's Lord hidden by the sentence of death (\work{In Matthaeum} 13.5, PL~9
col.~994), the Christ whom many took for the leaven according to Ambrose
(VII.187), Chrysologus's leaven of resurrection (\work{Sermo}~99) and
Chrysostom's disciples (\work{Hom.\ in Mt.}~46.2). So are the other readings
of the three measures, among them the Trinitarian reading on which Jerome
(the same place), Hilary (13.6, col.~995) and Chrysologus (\work{Sermo}~99) divide. Chrysostom
reads the verb ``hid'' of the disciples mixed among the nations, and the
speech says so; it uses his account of the verb and credits him with nothing
about love. The words ``the labour of love'' are the Greek that Chrysostom
expounds; the Missal's Latin, and the Douay--Rheims, read ``the work of your
faith and labour and charity'' (1~Thess 1:3; \latin{óperis fidei vestræ, et
labóris, et caritátis}), and the speech quotes that English and names the
Greek as his.

\textbf{Connections that are the homily's own.} These are drawn in the
preacher's voice and credited to no Father. Ambrose's teaching and
Augustine's love are said to come together in Bede's sentence, the commandment
being the teaching and love what it hides; the expansive study draws the same
inference (``Until the whole is leavened''). The remarks that hidden leaven
cannot be watched at work, that neither parable is about speed, and that a
part of us kept back from love stays as it was are the preacher's. The
Collect's meditating, always, is taken as our part in the long meditation
Ambrose ascribes to the Church, as in the expansive study's ``Thought, word
and deed''; Ambrose does not mention the Collect. The Collect's thought, word
and deed are read as the reach of the leaven through the whole person. That
Jason's deed is love gone past thought and word is the preacher's reading of
Chrysostom's example. The cry from the depths is said to work on the one who
prays as leaven works in the meal, the expansive study's inference from
Chrysostom on Ps 129:1, who does not mention the parable. Bede's meal and
Paul's one bread are joined to Communion by the preacher; Bede does not speak
of the Mass there. The response, to ask at Communion that love reach the whole
person and to believe that it will be received, applies the Communion's
promise (Mk 11:24) to the homily's argument; no Father cited comments on the
antiphon so. That those who have been fed ask to go on hungering because the
work is not finished in them reads the Postcommunion with Bede's ``begun here
and perfected there,'' as the expansive study does, and the closing judgement
on slowness is the preacher's conclusion.

\textbf{Sources, translations and loci.} Each source is listed once, with the
translation the speech uses; ``paraphrased'' means told in the speech's own
words. \emph{The Mass}, nos.~1759--1768. The Alleluia (Ps 129:1--2; Heb.\
130:1--2; p.~742) and the Communion (Mk 11:24; p.~743) are quoted in the
English of the Lasance \work{New Roman Missal} (1945), \work{Supplementary
Mass No.~1}, a historical study aid and not an approved liturgical
translation of the 1962 Missal; the Collect (p.~741) and the Postcommunion
(p.~743) are paraphrased in its words, with God named in the third person.
The Collect's \latin{semper rationabilia meditantes} is explained as meditating
always. \emph{Scripture} is quoted in the Douay--Rheims, Challoner revision,
at Mt 13:33 and 35, Mt 16:6, 1~Cor 5:8 and 1~Thess 1:3; Mt 13:1--2, the rest
of Mt 13:31--36 and Acts 17:1--10 are told close to the Douay's words.
\emph{St John Chrysostom}, in the Nicene and Post-Nicene English: \work{Hom.\ in Mt.}~47.1,
the disciples passing over the two parables ``as being plainer'',
paraphrased; 46.2, ``He said not, `put,' simply, but `hid','' paraphrased,
and ``though it be buried, yet is not destroyed, but by little and little
transmutes all into its own condition'', quoted; \work{Hom.\ in 1~Thess.}~1,
on 1:3, ``Merely to love is no labor at all. But to love genuinely is great
labor'' and ``when the seed had not yet taken root'', quoted, and the rest of
the passage on Jason paraphrased; \work{Expositio in Ps.}~129.1, paraphrased
from the Greek as the expansive study gives it. \emph{St Ambrose}, in the
expansive study's English of the Latin: \work{Expositio in Lucam} VII.190,
leaven as teaching shown from Mt 16:6 and 1~Cor 5:8, and \latin{aliud est
fermentum zizaniae, aliud fermentum tritici}, quoted; VII.190, the leaven of
the Church spiritual teaching, paraphrased; VII.191, \latin{non
tumultuarium, nec fortuitum, sed diuturna meditatione collatum; ut haec tria
unum sint}, paraphrased closely, with body, soul and spirit from
VII.190--191. \emph{St Augustine}, in the expansive study's English of the
Latin: \work{Quaestiones Evangeliorum} I.12, love \latin{quod fervefacit et
excitat}, paraphrased.
\emph{St Bede}, in the expansive study's English of the Migne Latin:
\work{In Lucam} IV, on Lk 13:19, \latin{crescit et in mente cujusque
credentis, quia nemo repente fit perfectus}; on Lk 13:21, the Church as the
woman and \latin{Abscondit ergo mulier fermentum dilectionis in farinae sata
tria, quia praecipit Ecclesia ut diligamus Dominum ex toto corde, tota anima,
tota virtute}; \latin{charitas in nostra mente recondita eousque crescere
debet, donec totam mentem in sui perfectionem commutet, ut nihil videlicet
praeter ipsam conditoris sui dilectionem diligere, agere, recolere anima
possit}; the meal \latin{quotquot timoris et spei exercitio, quasi mola
superiore et inferiore conterimur, ut juxta Apostolum unus panis unum corpus
multi simus in Christo} (1~Cor 10:17); and \latin{Quod hic quidem inchoari,
sed ibi habet perfici, ubi cum sit Deus omnia in omnibus, omnes uno eodemque
suae dilectionis igne calefecit} (1~Cor 15:28); all paraphrased closely.
The words for the meal, \latin{cujus farina nos sumus}, are a phrase Ambrose
had used (VII.187), and the love that warms and rouses is Augustine's gloss;
the speech names Augustine for it and gives Bede the millstones, the
commandment and the argument from ``until,'' which are his.
Bede is cited on Luke's form of the two parables, Ambrose on Luke's form of
the leaven.

\par}

\Needspace{14\baselineskip}
\section*{References}
\runninghead{References}

{\footnotesize
\begin{itemize}\setlength{\itemsep}{0.15em}\setlength{\parskip}{0pt}
\item \work{Missale Romanum ex decreto Sacrosancti Concilii Tridentini restitutum, Summorum Pontificum cura recognitum}, editio typica (Typis Polyglottis Vaticanis, 1962), CMAA facsimile, Proprium de Tempore, pp.~421--422, nos.~1759--1768.
\item Lasance, F.~X., and Walsh, F.~A., \work{The New Roman Missal in Latin and English}, revised edition (New York: Benziger Brothers, 1945), pp.~741--743 (\work{Supplementary Mass No.~1}: Collect, Alleluia, Communion, Postcommunion).
\item \work{The Holy Bible Translated from the Latin Vulgate} (Douay--Rheims, Challoner revision), Project Gutenberg text: Mt 13:1--2, 31--36; 16:6; Acts 17:1--10; 1~Cor 5:8; 1~Thess 1:3.
\item St Ambrose, \work{Expositio evangelii secundum Lucam} VII.187--191, Wikisource transcription of PL~15.
\item St Augustine, \work{Quaestiones Evangeliorum} I.12, Migne Latin text (Wikisource).
\item St Bede, \work{In Lucae evangelium expositio} IV, on Lk 13:19 and 13:21, PL~92, Wikisource transcription.
\item St Hilary of Poitiers, \work{In Evangelium Matthaei commentarius} 13.5--6, PL~9 (Paris, 1844), cols.~994--995.
\item St Jerome, \work{Commentarii in Matthaeum} II, on 13:33, PL~26 (Paris, 1845), cols.~91--92.
\item St John Chrysostom, \work{Homilies on Matthew} 46 and 47, Nicene and Post-Nicene Fathers, 1st ser., vol.~10; \work{Homilies on First Thessalonians} 1, Nicene and Post-Nicene Fathers, 1st ser., vol.~13; \work{Expositio in Psalmos}, on Ps 129, PG~55.
\item St Peter Chrysologus, \work{Sermo}~99, PL~52 (Paris, 1894).
\end{itemize}
\par}
